\documentclass[a4paper,12pt]{article}

\usepackage[T2A]{fontenc}
\usepackage[utf8]{inputenc}
\usepackage[russian]{babel}

\usepackage{amsmath,amssymb,amsthm,mathtools}
\usepackage{helvet}
\renewcommand{\familydefault}{\sfdefault}

\usepackage{icomma}
\usepackage[
  left=22mm,
  right=22mm,
  top=22mm,
  bottom=24mm,
  headheight=15pt
]{geometry}
\usepackage{microtype}
\usepackage{tikz}
\usepackage{tkz-euclide}
\usepackage{pgfplots}
\pgfplotsset{compat=1.18}
\usepackage{tabularx,booktabs,longtable}
\usepackage{enumitem}
\usepackage{hyperref}
\usepackage{parskip}
\usepackage{fancyhdr}
\usepackage{setspace}
\usepackage{xcolor}

\definecolor{accent}{HTML}{008080}
\definecolor{darkaccent}{HTML}{004D4D}
\definecolor{textgray}{HTML}{333333}
\definecolor{softgray}{HTML}{F2F5F5}
\definecolor{muted}{HTML}{6D7A7A}

\hypersetup{
  colorlinks=true,
  linkcolor=darkaccent,
  urlcolor=darkaccent,
  pdfauthor={Лёвин Артём Александрович},
  pdftitle={Парабола: построение графика и выколотые точки}
}

\onehalfspacing
\setlength{\parindent}{0pt}
\setlength{\emergencystretch}{2em}

\pagestyle{fancy}
\fancyhf{}
\fancyhead[L]{\small\color{textgray}Парабола и выколотые точки}
\fancyhead[R]{\small\color{textgray}София Кальней}
\fancyfoot[C]{\small\color{muted}\thepage}
\renewcommand{\headrulewidth}{0.4pt}
\renewcommand{\headrule}{\hbox to\headwidth{\color{accent}\leaders\hrule height \headrulewidth\hfill}}

\setlist[itemize]{
  leftmargin=1.6em,
  itemsep=0.3em,
  topsep=0.4em
}
\setlist[enumerate]{
  leftmargin=2em,
  itemsep=0.45em,
  topsep=0.4em
}

\newcommand{\term}[1]{\textcolor{darkaccent}{\textbf{#1}}}
\newcommand{\important}[1]{\textcolor{darkaccent}{#1}}
\newcommand{\checkitem}{\item[$\square$]}
\newcommand{\R}{\mathbb{R}}

\begin{document}

% ----------------------------------------------------------------
% ТИТУЛЬНЫЙ ЛИСТ
% ----------------------------------------------------------------

\begin{titlepage}
  \thispagestyle{empty}
  \centering

  \vspace*{2.2cm}

  {\Large\color{muted}\textbf{ЧЕК-ЛИСТ}\par}

  \vspace{1.1cm}

  {\fontsize{27}{33}\selectfont
   \color{darkaccent}\textbf{Парабола}\par}

  \vspace{0.35cm}

  {\Large\color{textgray}
   Построение графика функции\\[0.2cm]
   и работа с выколотыми точками\par}

  \vspace{1.7cm}

  \begin{tikzpicture}
    \draw[accent, line width=1.1pt]
      (-4.8,0) -- (4.8,0);
    \fill[accent] (0,0) circle (2pt);
  \end{tikzpicture}

  \vspace{1.7cm}

  {\large
   Лёвин Артём Александрович\\[0.2cm]
   эксклюзивно для Софии Кальней\par}

  \vspace{0.7cm}

  {\normalsize ОГЭ по математике\par}

  \vfill

  {\large\color{textgray}\today\par}

  \vspace{1.2cm}

  \begin{tikzpicture}[x=1cm,y=1cm]
    \draw[
      color=accent!55,
      line width=1.3pt,
      line cap=round
    ]
      (-7.3,0)
      .. controls (-4.8,0.75) and (-2.8,-0.55) .. (-0.6,0.05)
      .. controls (1.6,0.65) and (3.8,-0.65) .. (7.3,0.15);

    \draw[
      color=accent!25,
      line width=0.7pt,
      line cap=round
    ]
      (-6.7,-0.35)
      .. controls (-3.8,0.2) and (-1.8,-0.8) .. (0.6,-0.25)
      .. controls (3.0,0.3) and (4.8,-0.35) .. (6.8,-0.1);
  \end{tikzpicture}

\end{titlepage}

% ----------------------------------------------------------------
% ОСНОВНАЯ ЧАСТЬ
% ----------------------------------------------------------------

\section*{1. Квадратичная функция и парабола}

\term{Квадратичная функция} задаётся формулой
\[
  y=ax^2+bx+c,
  \qquad a\ne 0.
\]

Её график называется \term{параболой}.

В записи \(y=ax^2+bx+c\):

\begin{itemize}
  \item \(a\) --- коэффициент при \(x^2\);
  \item \(b\) --- коэффициент при \(x\);
  \item \(c\) --- свободный коэффициент;
  \item \(V(x_V;y_V)\) --- вершина параболы.
\end{itemize}

Через вершину проходит \term{ось симметрии} параболы:
\[
  x=x_V.
\]

\subsection*{Куда направлены ветви}

Направление ветвей определяется знаком коэффициента \(a\).

\begin{table}[ht]
  \centering
  \caption{Влияние коэффициента \(a\)}
  \renewcommand{\arraystretch}{1.35}
  \begin{tabularx}{0.88\linewidth}{@{}cX@{}}
    \toprule
    \textbf{Условие} & \textbf{Положение параболы} \\
    \midrule
    \(a>0\) & Ветви направлены вверх. Вершина является самой нижней точкой графика. \\
    \(a<0\) & Ветви направлены вниз. Вершина является самой верхней точкой графика. \\
    \bottomrule
  \end{tabularx}
\end{table}

Модуль коэффициента \(a\) влияет на форму графика:

\begin{itemize}
  \item при большом \(|a|\) ветви располагаются ближе к оси симметрии;
  \item при малом положительном \(|a|\) ветви располагаются шире.
\end{itemize}

\section*{2. Универсальный алгоритм построения параболы}

Для функции
\[
  y=ax^2+bx+c
\]
используйте следующий порядок действий.

\begin{enumerate}[label=\textcolor{darkaccent}{\textbf{\arabic*.}}]
  \item Определите знак коэффициента \(a\) и направление ветвей.
  \item Найдите координаты вершины.
  \item Найдите точки пересечения с осью \(Ox\).
  \item Найдите точку пересечения с осью \(Oy\).
  \item При необходимости вычислите дополнительные симметричные точки.
  \item Выберите масштаб и нанесите найденные точки.
  \item Проведите плавную кривую с учётом оси симметрии.
\end{enumerate}

\section*{3. Вершина параболы}

Абсцисса вершины вычисляется по формуле
\[
  \boxed{x_V=-\frac{b}{2a}}.
\]

Ординату вершины получают подстановкой найденного значения в исходную функцию:
\[
  \boxed{y_V=f(x_V)}.
\]

Таким образом,
\[
  V(x_V;y_V).
\]

Ось симметрии имеет уравнение
\[
  \boxed{x=x_V}.
\]

\subsection*{Проверка знаков}

В формуле
\[
  x_V=-\frac{b}{2a}
\]
первый минус относится ко всему коэффициенту \(b\).

Например, для функции
\[
  y=x^2-x-6
\]
имеем
\[
  a=1,\qquad b=-1,\qquad c=-6.
\]

Тогда
\[
  x_V=-\frac{-1}{2\cdot 1}=\frac12.
\]

Ордината вершины:
\[
  y_V=
  \left(\frac12\right)^2
  -\frac12-6
  =
  \frac14-\frac12-6
  =
  -\frac{25}{4}.
\]

Следовательно,
\[
  V\left(\frac12;-\frac{25}{4}\right).
\]

\section*{4. Пересечения с координатными осями}

\subsection*{Пересечение с осью \(Ox\)}

На оси \(Ox\) выполняется условие
\[
  y=0.
\]

Поэтому требуется решить квадратное уравнение
\[
  ax^2+bx+c=0.
\]

Если корни равны \(x_1\) и \(x_2\), точки пересечения имеют координаты
\[
  A(x_1;0),
  \qquad
  B(x_2;0).
\]

Количество точек пересечения определяется дискриминантом:
\[
  D=b^2-4ac.
\]

\begin{table}[ht]
  \centering
  \caption{Количество пересечений параболы с осью \(Ox\)}
  \renewcommand{\arraystretch}{1.35}
  \begin{tabularx}{0.8\linewidth}{@{}cX@{}}
    \toprule
    \textbf{Дискриминант} & \textbf{Количество точек} \\
    \midrule
    \(D>0\) & Две различные точки пересечения. \\
    \(D=0\) & Одна точка; парабола касается оси \(Ox\). \\
    \(D<0\) & Точки пересечения отсутствуют. \\
    \bottomrule
  \end{tabularx}
\end{table}

\subsection*{Пересечение с осью \(Oy\)}

На оси \(Oy\) выполняется условие
\[
  x=0.
\]

Подставим \(x=0\):
\[
  y=a\cdot 0^2+b\cdot 0+c=c.
\]

Следовательно, парабола всегда проходит через точку
\[
  \boxed{Q(0;c)}.
\]

Свободный коэффициент \(c\) является ординатой точки пересечения графика с осью \(Oy\).

\section*{5. Дополнительные точки}

Вершины и пересечений с осями часто достаточно для построения эскиза.

Если положение ветвей вызывает сомнения, составьте таблицу значений. Выбирайте абсциссы симметрично относительно \(x_V\):
\[
  x_V-d
  \qquad\text{и}\qquad
  x_V+d.
\]

Для таких значений функции совпадают:
\[
  f(x_V-d)=f(x_V+d).
\]

Например, при \(x_V=-2\) удобно взять пары
\[
  -3\ \text{и}\ -1,
  \qquad
  -4\ \text{и}\ 0.
\]

\section*{6. Полный пример построения параболы}

Построим график функции
\[
  y=x^2-x-6.
\]

\subsection*{Шаг 1. Направление ветвей}

\[
  a=1>0.
\]

Ветви направлены вверх.

\subsection*{Шаг 2. Вершина}

\[
  x_V=-\frac{-1}{2\cdot 1}=\frac12,
\]
\[
  y_V=
  \left(\frac12\right)^2
  -\frac12-6
  =
  -\frac{25}{4}.
\]

\[
  V\left(\frac12;-\frac{25}{4}\right).
\]

Ось симметрии:
\[
  x=\frac12.
\]

\subsection*{Шаг 3. Пересечение с осью \(Ox\)}

\[
  x^2-x-6=0,
\]
\[
  (x+2)(x-3)=0.
\]

Отсюда
\[
  x_1=-2,
  \qquad
  x_2=3.
\]

Точки:
\[
  A(-2;0),
  \qquad
  B(3;0).
\]

\subsection*{Шаг 4. Пересечение с осью \(Oy\)}

\[
  x=0,
  \qquad
  y=-6.
\]

Точка:
\[
  Q(0;-6).
\]

\begin{center}
\begin{tikzpicture}
  \begin{axis}[
    width=0.84\linewidth,
    height=8.6cm,
    axis lines=middle,
    unit vector ratio=1 1,
    xmin=-3.4,
    xmax=4.4,
    ymin=-7.2,
    ymax=8.2,
    samples=250,
    domain=-2.9:3.9,
    clip=false,
    grid=both,
    minor tick num=1,
    major grid style={line width=0.3pt,draw=gray!35},
    minor grid style={line width=0.2pt,draw=gray!18},
    xlabel={$x$},
    ylabel={$y$},
    tick label style={font=\small},
    label style={font=\small},
    every axis plot/.append style={line width=1.25pt}
  ]
    \addplot[accent] {x^2-x-6};

    \addplot[
      only marks,
      mark=*,
      mark size=2.4pt,
      darkaccent
    ] coordinates {
      (-2,0)
      (3,0)
      (0,-6)
      (0.5,-6.25)
    };

    \addplot[
      darkaccent,
      dashed,
      domain=-6.25:7
    ] ({0.5},x);

    \node[anchor=north east, font=\small]
      at (axis cs:-2,0) {$A(-2;0)$};

    \node[anchor=north west, font=\small]
      at (axis cs:3,0) {$B(3;0)$};

    \node[anchor=north east, font=\small]
      at (axis cs:0,-6) {$Q(0;-6)$};

    \node[anchor=north west, font=\small]
      at (axis cs:0.5,-6.25)
      {$V\left(\frac12;-\frac{25}{4}\right)$};
  \end{axis}
\end{tikzpicture}
\end{center}

\section*{7. Масштаб и оформление графика}

Перед построением оцените диапазон координат найденных точек.

\begin{itemize}
  \item Масштаб по осям \(Ox\) и \(Oy\) разрешается выбирать различным.
  \item Одинаковый масштаб обычно делает график более наглядным.
  \item Все ключевые точки должны помещаться на координатной плоскости.
  \item Координаты вершины, пересечений и выколотых точек следует подписывать.
  \item Ветви параболы проводят плавно и продолжают до границ выбранного окна.
\end{itemize}

При каждом фиксированном значении \(x\) функция имеет единственное значение \(y\). Вертикальная прямая должна пересекать график функции максимум в одной точке.

\section*{8. Важное различие при возведении в квадрат}

Записи
\[
  -x^2
  \qquad\text{и}\qquad
  (-x)^2
\]
имеют разный смысл.

\[
  -x^2=-(x^2),
\]
\[
  (-x)^2=x^2.
\]

Например,
\[
  -1{,}5^2=-2{,}25,
\]
поскольку степень относится только к числу \(1{,}5\).

При подстановке отрицательного значения переменной используйте скобки:
\[
  (-3)^2=9.
\]

Для функции
\[
  y=-x^2+3x+4
\]
при \(x=\frac32\):
\[
  y=
  -\left(\frac32\right)^2
  +3\cdot\frac32+4
  =
  -\frac94+\frac92+4
  =
  \frac{25}{4}.
\]

Следовательно,
\[
  V\left(\frac32;\frac{25}{4}\right).
\]

\section*{9. Графики алгебраических дробей}

Рассмотрим функцию
\[
  y=
  \frac{(x+4)(x^2+3x+2)}{x+1}.
\]

\subsection*{Шаг 1. Ограничение}

Знаменатель должен быть отличен от нуля:
\[
  x+1\ne 0,
\]
\[
  \boxed{x\ne -1}.
\]

Ограничение выписывается по исходному выражению и сохраняется до конца решения.

\subsection*{Шаг 2. Разложение на множители}

\[
  x^2+3x+2=(x+1)(x+2).
\]

Тогда
\[
  y=
  \frac{(x+4)(x+1)(x+2)}{x+1}.
\]

При \(x\ne -1\) получаем
\[
  y=(x+4)(x+2).
\]

Раскроем скобки:
\[
  y=x^2+6x+8,
  \qquad
  x\ne -1.
\]

График совпадает с параболой
\[
  y=x^2+6x+8
\]
во всех разрешённых точках. Точка с абсциссой \(x=-1\) исключена.

\subsection*{Шаг 3. Основные точки параболы}

Вершина:
\[
  x_V=-\frac{6}{2\cdot 1}=-3,
\]
\[
  y_V=(-3)^2+6\cdot(-3)+8=-1.
\]

\[
  V(-3;-1).
\]

Пересечения с осью \(Ox\):
\[
  x^2+6x+8=0,
\]
\[
  (x+4)(x+2)=0,
\]
\[
  x_1=-4,
  \qquad
  x_2=-2.
\]

Точки:
\[
  A(-4;0),
  \qquad
  B(-2;0).
\]

Пересечение с осью \(Oy\):
\[
  Q(0;8).
\]

\subsection*{Шаг 4. Координаты выколотой точки}

Подставим исключённое значение \(x=-1\) в сокращённую формулу:
\[
  y=(-1)^2+6\cdot(-1)+8=3.
\]

Получаем выколотую точку
\[
  \boxed{K(-1;3)}.
\]

Она отмечается незакрашенным кружком.

\begin{center}
\begin{tikzpicture}
  \begin{axis}[
    width=0.84\linewidth,
    height=8.6cm,
    axis lines=middle,
    unit vector ratio=1 1,
    xmin=-5.5,
    xmax=0.5,
    ymin=-2,
    ymax=10,
    samples=250,
    clip=false,
    grid=both,
    minor tick num=1,
    major grid style={line width=0.3pt,draw=gray!35},
    minor grid style={line width=0.2pt,draw=gray!18},
    xlabel={$x$},
    ylabel={$y$},
    tick label style={font=\small},
    label style={font=\small},
    every axis plot/.append style={line width=1.25pt}
  ]
    \addplot[
      accent,
      domain=-5.25:-1.03
    ] {x^2+6*x+8};

    \addplot[
      accent,
      domain=-0.97:0.2
    ] {x^2+6*x+8};

    \addplot[
      only marks,
      mark=*,
      mark size=2.4pt,
      darkaccent
    ] coordinates {
      (-3,-1)
      (-4,0)
      (-2,0)
      (0,8)
    };

    \addplot[
      only marks,
      mark=o,
      mark size=3.2pt,
      mark options={line width=1.2pt,draw=darkaccent,fill=white}
    ] coordinates {
      (-1,3)
    };

    \addplot[
      darkaccent,
      dashed,
      domain=-1:9
    ] ({-3},x);

    \node[anchor=north west, font=\small]
      at (axis cs:-3,-1) {$V(-3;-1)$};

    \node[anchor=south west, font=\small]
      at (axis cs:-1,3) {$K(-1;3)$};

    \node[anchor=south east, font=\small]
      at (axis cs:-4,0) {$A$};

    \node[anchor=south west, font=\small]
      at (axis cs:-2,0) {$B$};

    \node[anchor=south west, font=\small]
      at (axis cs:0,8) {$Q$};
  \end{axis}
\end{tikzpicture}
\end{center}

\section*{10. Алгоритм работы с выколотыми точками}

Для функции вида
\[
  y=\frac{P(x)}{Q(x)}
\]
используйте следующий порядок.

\begin{enumerate}[label=\textcolor{darkaccent}{\textbf{\arabic*.}}]
  \item Выпишите ограничения из исходного знаменателя:
  \[
    Q(x)\ne 0.
  \]

  \item Разложите числитель и знаменатель на множители.

  \item Сократите одинаковые множители.

  \item Сохраните все ограничения, найденные до сокращения.

  \item Постройте график упрощённой функции.

  \item Для каждого исключённого значения \(x_0\) вычислите
  \[
    y_0=f_{\text{упр}}(x_0).
  \]

  \item Отметьте точку \((x_0;y_0)\) незакрашенным кружком.
\end{enumerate}

\subsection*{Главное правило}

Сокращение общего множителя упрощает формулу, при этом исключённые значения переменной сохраняются:
\[
  \frac{(x-a)F(x)}{x-a}=F(x),
  \qquad x\ne a.
\]

В точке
\[
  \bigl(a;F(a)\bigr)
\]
на графике остаётся выкол.

\section*{11. Горизонтальная прямая \(y=m\)}

Прямая
\[
  y=m
\]
горизонтальна.

Количество общих точек прямой и графика равно количеству разрешённых решений уравнения
\[
  f(x)=m.
\]

Для обычной параболы горизонтальная прямая:

\begin{itemize}
  \item пересекает график в двух точках при расположении выше минимума параболы с ветвями вверх;
  \item касается графика в вершине;
  \item общих точек не имеет при расположении ниже вершины параболы с ветвями вверх.
\end{itemize}

Для параболы с ветвями вниз положения меняются симметрично.

Выколотая точка способна уменьшить количество общих точек. Если одна из двух точек пересечения исключена, остаётся ровно одна общая точка.

\section*{12. Контрольный чек-лист построения}

Перед завершением решения проверьте каждый пункт.

\begin{itemize}
  \checkitem Коэффициенты \(a\), \(b\), \(c\) выписаны верно.
  \checkitem Знак коэффициента \(a\) учтён.
  \checkitem Формула \(x_V=-\dfrac{b}{2a}\) применена со всеми знаками.
  \checkitem Значение \(y_V\) найдено подстановкой в исходную функцию.
  \checkitem Ось симметрии \(x=x_V\) отмечена.
  \checkitem Для пересечения с \(Ox\) использовано условие \(y=0\).
  \checkitem Для пересечения с \(Oy\) использовано условие \(x=0\).
  \checkitem Ограничения выписаны до сокращения дроби.
  \checkitem Координаты выколотых точек вычислены по упрощённой формуле.
  \checkitem Выколотые точки обозначены незакрашенными кружками.
  \checkitem Масштаб позволяет увидеть все ключевые точки.
  \checkitem Парабола проведена плавно и симметрично.
  \checkitem Координаты основных точек подписаны.
\end{itemize}

\section*{13. Типичные ошибки}

\begin{enumerate}
  \item Потеря первого минуса в формуле
  \[
    x_V=-\frac{b}{2a}.
  \]

  \item Подстановка отрицательного числа без скобок:
  \[
    (-3)^2=9.
  \]

  \item Смешение условий для координатных осей:
  \[
    Ox:\ y=0,
    \qquad
    Oy:\ x=0.
  \]

  \item Возврат исключённого значения после сокращения дроби.

  \item Выкалывание точки без вычисления её ординаты.

  \item Неверное положение выколотой точки из-за приблизительного чтения графика.

  \item Построение ломаной линии вместо плавной параболы.

  \item Загиб ветви, при котором одному значению \(x\) соответствуют два значения \(y\).

  \item Слишком мелкий масштаб, скрывающий вершину или точки пересечения.

  \item Отсутствие подписей у ключевых точек.
\end{enumerate}

% ----------------------------------------------------------------
% ТРЕНИРОВОЧНЫЙ БЛОК
% ----------------------------------------------------------------

\clearpage
\section*{Тренировочный блок}

Во всех заданиях выполните аккуратный эскиз графика. Отметьте вершину, ось симметрии, пересечения с координатными осями и все выколотые точки.

\subsection*{Средний уровень}

\begin{enumerate}[label=\textbf{\arabic*.}]
  \item Постройте график функции
  \[
    y=x^2-4x+3.
  \]

  \item Постройте график функции
  \[
    y=-x^2+2x+3.
  \]

  \item Постройте график функции
  \[
    y=2x^2+8x+6.
  \]
\end{enumerate}

\subsection*{Уровень выше среднего}

\begin{enumerate}[label=\textbf{\arabic*.},start=4]
  \item Постройте график функции
  \[
    y=
    \frac{(x-1)(x+2)(x-3)}{x-1}.
  \]

  \item Постройте график функции
  \[
    y=
    \frac{(x+4)(x^2+3x+2)}{x+1}.
  \]

  \item Постройте график функции
  \[
    y=
    \frac{(x-2)(x^2-5x+6)}{x-3}.
  \]

  \item Постройте график функции
  \[
    y=
    \frac{(x+1)(x-4)(-x^2+2x+8)}
         {(x+1)(x-4)}.
  \]

  \item Постройте график функции
  \[
    y=
    \frac{x(x^2-4x+3)}{x}.
  \]
  Определите все значения \(m\), при которых прямая
  \[
    y=m
  \]
  имеет с графиком ровно одну общую точку.

  \item Постройте график функции
  \[
    y=
    \frac{(x-1)(-x^2+6x-5)}{x-1}.
  \]
  Определите все значения \(m\), при которых прямая
  \[
    y=m
  \]
  имеет с графиком ровно одну общую точку.
\end{enumerate}

\subsection*{Сложный уровень}

\begin{enumerate}[label=\textbf{\arabic*.},start=10]
  \item Постройте график функции
  \[
    y=
    \frac{x(x+1)^2(x-3)}
         {x(x+1)}.
  \]
  Определите все значения \(m\), при которых прямая
  \[
    y=m
  \]
  имеет с графиком ровно одну общую точку.
\end{enumerate}

% ----------------------------------------------------------------
% ОТВЕТЫ
% ----------------------------------------------------------------

\clearpage
\section*{Ответы к тренировочному блоку}

\small
\renewcommand{\arraystretch}{1.35}

\begin{longtable}{@{}p{0.07\linewidth}p{0.87\linewidth}@{}}
  \caption{Ключевые данные для проверки построений}\\
  \toprule
  \textbf{№} & \textbf{Ответ} \\
  \midrule
  \endfirsthead

  \toprule
  \textbf{№} & \textbf{Ответ} \\
  \midrule
  \endhead

  \bottomrule
  \endfoot

  1 &
  Ветви вверх.
  \[
    V(2;-1),\qquad x=2.
  \]
  Пересечения:
  \[
    (1;0),\quad(3;0),\quad(0;3).
  \]
  \\

  2 &
  Ветви вниз.
  \[
    V(1;4),\qquad x=1.
  \]
  Пересечения:
  \[
    (-1;0),\quad(3;0),\quad(0;3).
  \]
  \\

  3 &
  Ветви вверх.
  \[
    V(-2;-2),\qquad x=-2.
  \]
  Пересечения:
  \[
    (-3;0),\quad(-1;0),\quad(0;6).
  \]
  \\

  4 &
  Ограничение:
  \[
    x\ne 1.
  \]
  После сокращения:
  \[
    y=x^2-x-6.
  \]
  \[
    V\left(\frac12;-\frac{25}{4}\right).
  \]
  Корни:
  \[
    -2,\quad 3.
  \]
  Пересечение с \(Oy\):
  \[
    (0;-6).
  \]
  Выколотая точка:
  \[
    (1;-6).
  \]
  \\

  5 &
  Ограничение:
  \[
    x\ne -1.
  \]
  После сокращения:
  \[
    y=x^2+6x+8.
  \]
  \[
    V(-3;-1).
  \]
  Пересечения:
  \[
    (-4;0),\quad(-2;0),\quad(0;8).
  \]
  Выколотая точка:
  \[
    (-1;3).
  \]
  \\

  6 &
  Ограничение:
  \[
    x\ne 3.
  \]
  После сокращения:
  \[
    y=(x-2)^2.
  \]
  \[
    V(2;0),\qquad x=2.
  \]
  Пересечения:
  \[
    (2;0),\quad(0;4).
  \]
  Выколотая точка:
  \[
    (3;1).
  \]
  \\

  7 &
  Ограничения:
  \[
    x\ne -1,\qquad x\ne 4.
  \]
  После сокращения:
  \[
    y=-x^2+2x+8.
  \]
  \[
    V(1;9).
  \]
  Разрешённое пересечение с \(Ox\):
  \[
    (-2;0).
  \]
  Пересечение с \(Oy\):
  \[
    (0;8).
  \]
  Выколотые точки:
  \[
    (-1;5),\qquad(4;0).
  \]
  \\

  8 &
  Ограничение:
  \[
    x\ne 0.
  \]
  После сокращения:
  \[
    y=x^2-4x+3.
  \]
  \[
    V(2;-1).
  \]
  Выколотая точка:
  \[
    (0;3).
  \]
  Ровно одна общая точка получается при
  \[
    \boxed{m=-1\quad\text{или}\quad m=3}.
  \]
  \\

  9 &
  Ограничение:
  \[
    x\ne 1.
  \]
  После сокращения:
  \[
    y=-x^2+6x-5.
  \]
  \[
    V(3;4).
  \]
  Выколотая точка:
  \[
    (1;0).
  \]
  Ровно одна общая точка получается при
  \[
    \boxed{m=0\quad\text{или}\quad m=4}.
  \]
  \\

  10 &
  Ограничения:
  \[
    x\ne -1,\qquad x\ne 0.
  \]
  После сокращения:
  \[
    y=(x+1)(x-3)=x^2-2x-3.
  \]
  \[
    V(1;-4).
  \]
  Выколотые точки:
  \[
    (-1;0),\qquad(0;-3).
  \]
  Ровно одна общая точка получается при
  \[
    \boxed{m=-4,\quad m=-3,\quad m=0}.
  \]
  \\

\end{longtable}

\normalsize

\end{document}